% ===== ch09 : 原文 main.tex 第 2671--3290 行（5d 限制到子群概形 / 5e 群概形作用的乘积）=====
% 本文件是原文 LaTeX 源码的逐行分片（去掉整段注释掉的旧稿），仅供翻译时对照；最终译文写到 src/chunks/ch09.tex。
\subsection{Restriction to subgroup schemes.}

Let $G$ be a group scheme,
and $H\subseteq G$ a closed subgroup scheme
with the inclusion morphism $i_H : H \hookrightarrow G$.
A $G$-action
$\sigma_G : G \times X \rightarrow X$
on a scheme $X$ naturally induces a $H$-action on $X$ defined by
$$
% ---- 原文行 2680 ----
\sigma_H
:=
(i_H \times \id_X)
\comp
\sigma
:
H\times X\rightarrow X.
$$


\begin{defn}[Restriction of equivariant sheaves]\label{defn:RestrictionOfEquivariantSheaves}
For any $G$-equivariant $\mathcal{O}_X$-module $(\mathcal{F}, \theta) \in \Mod_G(X)$, its \emph{restriction} to $H$ is defined to be the $H$-equivariant $\mathcal{O}_X$-module
$$
(\mathcal{F}, \theta)|_H := (\mathcal{F}, \theta|_H),
$$
where the restricted $H$-linearization
is defined by the pullback of the $G$-linearization
$$
\theta|_H 
:= 
% ---- 原文行 2700 ----
(i_H \times \id_X)^* \theta 
: 
\sigma_H^* \mathcal{F} 
\xrightarrow{\cong} 
\left(\pr_X^{H \times X}\right)^* \mathcal{F}
$$

This operation on objects, together with the identity map on morphisms, naturally defines a functor
$$
\Res_H^G : \Mod_G(X) \rightarrow \Mod_H(X),
$$
called the \emph{restriction functor}.
This also induces restriction functors on the corresponding full subcategories:
$$
\QCoh_G(X) \rightarrow \QCoh_H(X),
\quad
\Coh_G(X) \rightarrow \Coh_H(X),
\quad\text{and}\quad
\VB_G(X) \rightarrow \VB_H(X).
$$
% ---- 原文行 2720 ----
\end{defn}


The following lemma is standard:

\begin{lemma}\cite[Exercise 2.8]{Waterhouse}
\label{lem:DiagonalizableSurjective}
Let $G$ be a diagonalizable group scheme
over a field $\field$,
and let $H \subseteq G$ be a closed subgroup scheme.
Then $H$ is also diagonalizable, and the restriction map on character groups
$
i_H^* : \widehat{G} \rightarrow \widehat{H}
$
is surjective with $\ker(i_H^*) = H^\perp$,
where
$$
H^\perp
:=
\{ \cha \in \widehat{G} \mid \cha(h) = 1 \text{ for all } h \in H \}
% ---- 原文行 2740 ----
\subseteq
\widehat{G}
$$
is the annihilator of $H$ in $\widehat{G}$.
\end{lemma}

\begin{proof}
Since $H$ is a closed subgroup scheme of
a diagonalizable group scheme
$G\cong \Spec(\field[\widehat{G}])$,
there is a Hopf ideal $I$ of $\field[\widehat{G}]$
such that
$H\cong\Spec(\field[\widehat{G}]/I)$
and
$\mathcal{O}_H(H)\cong\field[\widehat{G}]/I$.

We claim that
$$
I=\left<\cha-1|\cha\in H^\perp\right>
\subseteq
% ---- 原文行 2760 ----
\field[\widehat{G}].
$$

For any
$\cha \in H^\perp\subseteq\field[\widehat{G}]\cong\mathcal{O}_G(G)$,
we have
$\cha|_H=1_{\mathcal{O}_H(H)}\in\mathcal{O}_H(H)$,
and hence
$(\cha-1)|_H=0\in \mathcal{O}_H(H)$.
This implies
$\langle \cha - 1 \mid \cha \in H^\perp \rangle \subseteq I$.

Conversely,
suppose
$f = \sum_{\cha \in \widehat{G}} c_{\cha} \cha \in I\subseteq\field[\widehat{G}]\cong\mathcal{O}_G(G)$.
Grouping the terms by their restricted images in $\mathcal{O}_H(H)$, we evaluate $f|_H$:
$$
0 = f|_H = \sum_{\mu} \left( \sum_{\cha|_H = \mu} c_{\cha} \right) \mu \quad \text{in } \mathcal{O}_H(H),
$$
where $\mu$ ranges over the distinct non-zero images of characters in $\mathcal{O}_H(H)$. Since each $\mu$ is a group-like element in the Hopf algebra $\mathcal{O}_H(H)$, these distinct elements $\mu$ are linearly independent over $\field$
% ---- 原文行 2780 ----
\cite[\S2.2]{Waterhouse}.
Consequently, the coefficient of each $\mu$ must vanish; that is, $\sum_{\cha|_H = \mu} c_{\cha} = 0$ for all such $\mu$. 

For each $\mu$ appearing in the sum, fix a representative $\cha_0 \in \widehat{G}$ such that $\cha_0|_H = \mu$. Any other character $\cha$ restricting to $\mu$ satisfies $(\cha \cha_0^{-1})|_H = 1$, meaning $\cha \cha_0^{-1} \in H^\perp$. Thus, we can write $\cha = h\cha_0$ for some $h \in H^\perp$. We can then rewrite the sub-sum corresponding to $\mu$ as follows:
$$
\sum_{h \in H^\perp} c_{h\cha_0} (h\cha_0) = \sum_{h \in H^\perp} c_{h\cha_0} (h - 1)\cha_0 + \left( \sum_{h \in H^\perp} c_{h\cha_0} \right) \cha_0.
$$
Since the sum of the coefficients $\sum c_{h\cha_0}$ is exactly zero, the second term vanishes. This leaves only a linear combination of elements of the form $h-1$ for $h \in H^\perp$. Summing over all such $\mu$, we conclude that $f \in \langle \cha - 1 \mid \cha \in H^\perp \rangle$.

Thus, $I = \langle \cha - 1 \mid \cha \in H^\perp \rangle$. By the first isomorphism theorem for rings, this yields the natural isomorphism $\mathcal{O}_H(H) \cong \field[\widehat{G}]/I \cong \field[\widehat{G}/H^\perp]$,
concluding that $H$ is diagonalizable
with
$\widehat{G}/H^\perp\cong\widehat{H}$.
\end{proof}
% \begin{proof}
% Since $H$ is a closed subgroup scheme of
% a diagonalizable group scheme
% $G\cong \Spec(\field[\widehat{G}])$,
% there is a Hopf ideal $I$ of $\field[\widehat{G}]$
% such that
% ---- 原文行 2800 ----
% $H\cong\Spec(\field[\widehat{G}]/I)$
% and
% $\mathcal{O}_H(H)\cong\field[\widehat{G}]/I$.

% Note that the natural inclusion
% $\widehat{H} \hookrightarrow \mathcal{O}_H(H)$
% extends by $\field$-linearity to an algebra homomorphism
% $$
% \phi:\field[\widehat{H}] \rightarrow \mathcal{O}_H(H),
% \quad
% \sum_{\cha\in\widehat{H}} c_{\cha}\cdot\cha \mapsto \sum_{\cha\in\widehat{H}} c_{\cha}\cha.
% $$

% On the other hand,
% the
% restriction map
% $
% i_H^* : \widehat{G} \rightarrow \widehat{H},
% $
% $
% ---- 原文行 2820 ----
% \cha\mapsto\cha|_H
% $
% induces a map of group algebras
% $$
% \field[i_H^*]:\field[\widehat{G}]\rightarrow\field[\widehat{H}],
% \quad
% \sum_{\cha\in\widehat{G}}c_{\cha}\cdot\cha
% \mapsto
% \sum_{\cha\in\widehat{G}}c_{\cha}\cdot\cha|_H.
% $$
% It factors through
% $\field[\widehat{G}]/I$,
% because
% $
% \sum_{\cha\in\widehat{G}}c_{\cha}\cdot\cha\in I\subseteq\mathcal{O}_G(G)
% $
% implies that it vanishes when it is restricted to $H$.
% Therefore,
% we have a natural map
% $$
% ---- 原文行 2840 ----
% \psi:\mathcal{O}_H(H)\rightarrow\field[\widehat{H}].
% $$

% Consider the composition
% $
% \psi\comp\phi
% :
% \mathcal{O}_H(H)\rightarrow\field[\widehat{H}]
% \rightarrow\mathcal{O}_H(H).
% $
% For any $\cha \in \widehat{G}$,
% its image $\overline{\cha}$ in
% $\mathcal{O}_H(H) \cong \field[\widehat{G}]/I$
% is mapped by $\psi$ to $\cha|_H \in \field[\widehat{H}]$,
% which $\phi$ then maps to $\cha|_H \in \mathcal{O}_H(H)$.
% Since $\cha|_H$ is exactly the equivalence class $\overline{\cha}$ modulo $I$,
% we have
% $\phi(\psi(\overline{\cha})) = \overline{\cha}$.
% Since the images of $\cha \in \widehat{G}$ span $\mathcal{O}_H(H)$,
% the composition $\psi \comp \phi$ is the identity map on $\mathcal{O}_H(H)$,
% ---- 原文行 2860 ----
% implying that $\phi$ is surjective.


% Furthermore, since distinct characters in $\widehat{H}$ are non-zero group-like elements in $\mathcal{O}_H(H)$, they are linearly independent over $\field$
% \cite[\S2.2]{Waterhouse}.
% This guarantees that $\phi$ is also injective. 
% Therefore, it is an isomorphism,
% concluding that $H$ is diagonalizable.



% As $\psi$ is surjective,
% the map
% $
% \field[i_H^*]
% :
% \field[\widehat{G}]\rightarrow\field[\widehat{H}]
% $
% is also surjective.
% This forces the underlying group homomorphism
% ---- 原文行 2880 ----
% $i_H^* : \widehat{G} \rightarrow \widehat{H}$
% to be surjective.
% \end{proof}

\begin{prop}[$H$-invariants of hom spaces]
\label{prop:HInvariantsOfHomSpace}
Let $G$ be a diagonalizable group scheme acting on a scheme $X$, and $H\subseteq G$ a closed subgroup scheme.
For any
$(\mathcal{F}, \theta)$,
$(\mathcal{F}', \theta') \in \Mod_G(X)$
such that
$\mathcal{F}$
is finitely presented,
the hom space
from
$(\mathcal{F}, \theta|_H)$
to
$(\mathcal{F}', \theta'|_H)$
in $\Mod_H(X)$
admits a canonical decomposition
% ---- 原文行 2900 ----
$$
\operatorname{Hom}_{\Mod_H(X)}\left( (\mathcal{F}, \theta|_H), (\mathcal{F}', \theta'|_H) \right)
=
\bigoplus_{\cha \in H^\perp} \operatorname{Hom}_{\Mod_G(X)}\left( (\mathcal{F}, \theta)(\cha), (\mathcal{F}', \theta') \right).
$$
\end{prop}




\begin{proof}
Note that $H$ is diagonalizable
by Lemma \ref{lem:DiagonalizableSurjective}.
Therefore,
by Proposition \ref{prop:WeightSpaceDecompositionOfHomSpace},
we have two decompositions
\begin{align*}
\operatorname{Hom}_{\Mod(X)}(\mathcal{F}, \mathcal{F}') 
&=
\bigoplus_{\cha \in \widehat{G}} \operatorname{Hom}_{\Mod_G(X)}\left( (\mathcal{F}, \theta)(\cha), (\mathcal{F}', \theta') \right)
% ---- 原文行 2920 ----
\\
&=
\bigoplus_{\chaa \in \widehat{H}} \operatorname{Hom}_{\Mod_H(X)}\left( (\mathcal{F}, \theta|_H)(\chaa), (\mathcal{F}', \theta'|_H) \right).
\end{align*}



We claim that
the following holds for each $\cha\in\widehat{G}$:
$$
\operatorname{Hom}_{\Mod_G(X)}\left( (\mathcal{F}, \theta)(\cha), (\mathcal{F}', \theta') \right)
\subseteq
\operatorname{Hom}_{\Mod_H(X)}\left( (\mathcal{F}, \theta|_H)(i_H^*\cha), (\mathcal{F}', \theta'|_H) \right).
$$
Indeed,
$
f\in\operatorname{Hom}_{\Mod_G(X)}\left( (\mathcal{F}, \theta)(\cha), (\mathcal{F}', \theta') \right)
$
satisfies
$$
% ---- 原文行 2940 ----
\sigma_G^* f\comp \theta'
=
\left(\left(\pr_G^{G\times X}\right)^*\cha \cdot \theta\right)
\comp
\left(\pr_X^{G\times X}\right)^* f
$$
Taking $(i_H^*\times \id_X)^*$
on both sides yields
$$
\sigma_H^*f\comp\left(\theta'|_H\right)
=
\left(
\left(\pr_H^{H\times X}\right)^*
i_H^*\cha
\cdot
\theta|_H
\right)
\comp
\left(\pr_X^{H\times X}\right)^*f,
$$
% ---- 原文行 2960 ----
where
$
\theta|_H
(i_H^*\cha)
=
\left(
\left(\pr_H^{H\times X}\right)^*
i_H^*\cha
\cdot
\theta|_H
\right).
$
This implies
$
f\in
\operatorname{Hom}_{\Mod_H(X)}\left( (\mathcal{F}, \theta|_H)(i_H^*\cha), (\mathcal{F}', \theta'|_H) \right),
$
establishing the claim.


% ---- 原文行 2980 ----
Therefore,
we have the following sub-decomposition
for each $\chaa\in\widehat{H}$:
$$
\operatorname{Hom}_{\Mod_H(X)}\left( (\mathcal{F}, \theta|_H)(\chaa), (\mathcal{F}', \theta'|_H) \right)
=
\bigoplus_{\substack{\cha \in \widehat{G} \\[1mm] i_H^*\cha = \chaa}} \operatorname{Hom}_{\Mod_G(X)}\left( (\mathcal{F}, \theta)(\cha), (\mathcal{F}', \theta') \right).
$$
Now the proposition is a special case when $\chaa=1$.
\end{proof}
Similar to corollary \ref{thm:EmbeddingOfTheOrbitCategory}, we have an embedding of orbit categories to restrict $G$-equivariant vector bundles to $H$-equivariant ones.
\begin{crl}\label{crl:orbit-category-GH-restriction}
    Let $G$ be a group scheme and $H\subseteq G$ be a subgroup scheme of $G$. Then there is an embedding of abelian categories \[
        \QCoh_G (X)/H^{\perp}\hookrightarrow\QCoh_H (X).
    \]
\end{crl}
\begin{proof}
By definition, the hom space in the orbit category is given by
$$
\operatorname{Hom}_{\QCoh_G(X)/H^\perp}\left( (\mathcal{F}, \theta), (\mathcal{F}', \theta') \right)
% ---- 原文行 3000 ----
=
\bigoplus_{\cha \in H^\perp} \operatorname{Hom}_{\QCoh_G(X)}\left( (\mathcal{F}, \theta)(\cha), (\mathcal{F}', \theta') \right).
$$
By Proposition \ref{prop:HInvariantsOfHomSpace},
this direct sum exactly coincides with the hom space
\[\operatorname{Hom}_{\QCoh_H(X)}
\left((\mathcal{F}, \theta|_H), (\mathcal{F}', \theta'|_H) \right)\]
in
$\QCoh_H(X)\subseteq\Mod_H(X)$.
Therefore, the functor $\QCoh_G(X)/H^\perp \hookrightarrow \QCoh_H(X)$ defined by $(\mathcal{F}, \theta) \mapsto (\mathcal{F},\theta|_H)$
on objects and by the identity map on morphisms is fully faithful.
\end{proof}

\subsection{Product of group scheme actions.}


Suppose that group schemes $G$ and $H$
act on schemes $X$ and $Y$
with action homomorphisms
$\sigma_X : G \times X \rightarrow X$
% ---- 原文行 3020 ----
and
$\sigma_Y : H \times Y \rightarrow Y$,
respectively.
They naturally induce an action of
the product group scheme $G\times H$
on the product scheme $X\times Y$
given by the action morphism
$$
\sigma_{X\times Y}
:=
\sigma_X\times \sigma_Y
:
(G\times H)\times (X\times Y)
\rightarrow
X\times Y.
$$
\begin{defn}[External tensor product of equivariant sheaves]
\label{defn:ExternalTensorProductOfEquivariantModules}
%Let $G$ and $H$ be group schemes acting on schemes $X$ and $Y$,
%respectively.
% ---- 原文行 3040 ----
Let $(\mathcal{F},\theta)\in\Mod_G(X)$ and
$(\mathcal{G},\eta)\in\Mod_H(Y)$.
The \emph{external tensor product} of
$(\mathcal{F},\theta)$ and $(\mathcal{G},\eta)$
$$
(\mathcal{F},\theta)\boxtimes(\mathcal{G},\eta)
:=
\left(
\mathcal{F}\boxtimes\mathcal{G},
\
\theta\boxtimes\eta
\right)
\in\Mod_{G\times H}(X\times Y)
$$
is defined as the
$\mathcal{O}_{X\times Y}$-module
$
\mathcal{F}\boxtimes\mathcal{G}
:=
\pi_X^*\mathcal{F}
% ---- 原文行 3060 ----
\underset{\mathcal{O}_{X\times Y}}{\otimes}
\pi_Y^*\mathcal{G}
$
together with the
$G\times H$-linearization
$$
\theta\boxtimes\eta
:=
\tilde{\pi}_{G\times X}^*\theta\otimes\tilde{\pi}_{H\times Y}^*\eta
:
\sigma_{X\times Y}^*(\mathcal{F}\boxtimes\mathcal{G})
\xrightarrow{\cong}
\left(\pr_{X\times Y}^{(G\times H)\times(X\times Y)}\right)^*
(\mathcal{F}\boxtimes\mathcal{G})
$$
where
\begin{align*}
\tilde{\pi}_{G\times X}^*\theta
&:
\tilde{\pi}_{G\times X}^*(\sigma_X^*\mathcal{F})
% ---- 原文行 3080 ----
\xrightarrow{\cong}
\tilde{\pi}_{G\times X}^*\left(\left(\pr_X^{G\times X}\right)^{\ast}\mathcal{F}\right),
\\[2mm]
\tilde{\pi}_{H\times Y}^*\eta
&:
\tilde{\pi}_{H\times Y}^*(\sigma_Y^*\mathcal{G})
\xrightarrow{\cong}
\tilde{\pi}_{H\times Y}^*\left(\left(\pr_Y^{H\times Y}\right)^*\mathcal{G}\right)
\end{align*}
are pull-backs of $\theta$ and $\eta$
by the natural projections
\begin{align*}
\tilde{\pi}_{G\times X}
&:=
\pr_{G\times X}^{(G\times H)\times(X\times Y)}
:
(G\times H)\times(X\times Y)
\rightarrow
G\times X,
\\[2mm]
% ---- 原文行 3100 ----
\tilde{\pi}_{H\times Y}
&:=
\pr_{H\times Y}^{(G\times H)\times(X\times Y)}
:
(G\times H)\times(X\times Y)
\rightarrow
H\times Y,
\end{align*}
respectively,
and
the tensor product
$
\tilde{\pi}_{G\times X}^*\theta\otimes\tilde{\pi}_{H\times Y}^*\eta
$
is given under the canonical isomorphisms
\begin{align*}
\sigma_{X\times Y}^*(\mathcal{F}\boxtimes\mathcal{G})
&=
\sigma_{X\times Y}^*
\Big(
% ---- 原文行 3120 ----
\pi_X^*\mathcal{F}
\underset{\mathcal{O}_{X\times Y}}{\otimes}
\pi_Y^*\mathcal{G}
\Big)
\\[2mm]
&\cong
(\sigma_{X\times Y}\comp\pi_X)^*\mathcal{F}
\underset{\mathcal{O}_{(G\times H)\times(X\times Y)}}{\otimes}
(\sigma_{X\times Y}\comp\pi_Y)^*\mathcal{G}
\\[2mm]
&=
(\tilde{\pi}_{G\times X}\comp \sigma_X)^*\mathcal{F}
\underset{\mathcal{O}_{(G\times H)\times(X\times Y)}}{\otimes}
(\tilde{\pi}_{H\times Y}\comp \sigma_Y)^*\mathcal{G}
\\[2mm]
&\cong
\tilde{\pi}_{G\times X}^*\left(\sigma_X^*\mathcal{F}\right)
\underset{\mathcal{O}_{(G\times H)\times(X\times Y)}}{\otimes}
\tilde{\pi}_{H\times Y}^*\left(\sigma_Y^*\mathcal{G}\right)
\end{align*}
% ---- 原文行 3140 ----
and
\begin{align*}
\left(\pr_{X\times Y}^{(G\times H)\times(X\times Y)}\right)^*
(\mathcal{F}\boxtimes\mathcal{G})
&=
\left(\pr_{X\times Y}^{(G\times H)\times(X\times Y)}\right)^*
\Big(
\pi_X^*\mathcal{F}
\underset{\mathcal{O}_{X\times Y}}{\otimes}
\pi_Y^*\mathcal{G}
\Big)
\\
&\cong
\left(\pr_{X\times Y}^{(G\times H)\times(X\times Y)}\comp\pi_X\right)^*\mathcal{F}
\underset{\mathcal{O}_{(G\times H)\times(X\times Y)}}{\otimes}
\left(\pr_{X\times Y}^{(G\times H)\times(X\times Y)}\comp\pi_Y\right)^*\mathcal{G}
\\
&\cong
\left(\tilde{\pi}_{G\times X}\comp\pr_X^{G\times X}\right)^*\mathcal{F}
\underset{\mathcal{O}_{(G\times H)\times(X\times Y)}}{\otimes}
% ---- 原文行 3160 ----
\left(\tilde{\pi}_{H\times Y}\comp\pr_Y^{H\times Y}\right)^*\mathcal{G}
\\
&\cong
\tilde{\pi}_{G\times X}^*\left(\left(\pr_X^{G\times X}\right)^*\mathcal{F}\right)
\underset{\mathcal{O}_{(G\times H)\times(X\times Y)}}{\otimes}
\tilde{\pi}_{H\times Y}^*\left(\left(\pr_Y^{H\times Y}\right)^*\mathcal{G}\right).
\end{align*}







For morphisms
$
f
\in
\Hom_{\Mod_G(X)}
((\mathcal{F},\theta),(\mathcal{F}',\theta'))
% ---- 原文行 3180 ----
$
and
$
g
\in
\Hom_{\Mod_H(Y)}
((\mathcal{G},\eta),(\mathcal{G}',\eta')),
$
define
$$
f\boxtimes g
:=
\pi_X^*f
\otimes
%\underset{\mathcal{O}_{X\times Y}}{\otimes}
\pi_Y^*g
\in
\Hom_{\Mod_{G\times H}(X\times Y)}
\big(
\left(
% ---- 原文行 3200 ----
\mathcal{F}\boxtimes\mathcal{G},
\
\theta\boxtimes\eta
\right),
\left(
\mathcal{F}'\boxtimes\mathcal{G}',
\
\theta'\boxtimes\eta'
\right)
\big).
$$
%The morphism $f\boxtimes g$ is $(G\times H)$-equivariant,
%since its equivariance condition is obtained by pulling back
%and tensoring the equivariance conditions for $f$ and $g$.

%$$
%f\boxtimes g
%:=
%\pi_X^*f
%\otimes
% ---- 原文行 3220 ----
%%\underset{\mathcal{O}_{X\times Y}}{\otimes}
%\pi_Y^*g
%:
%\mathcal{F}\boxtimes\mathcal{G}
%=
%\pi_X^*\mathcal{F}
%\underset{\mathcal{O}_{X\times Y}}{\otimes}
%\pi_Y^*\mathcal{G}
%\
%\rightarrow
%\
%\mathcal{F}'\boxtimes\mathcal{G}'
%=
%\pi_X^*\mathcal{F}'
%\underset{\mathcal{O}_{X\times Y}}{\otimes}
%\pi_Y^*\mathcal{G}'
%$$


This defines a bifunctor
% ---- 原文行 3240 ----
$$
\boxtimes:
\Mod_G(X)\times\Mod_H(Y)
\rightarrow
\Mod_{G\times H}(X\times Y),
$$
which restricts to equivariant (quasi-)coherent sheaves and
equivariant vector bundles.
\end{defn}



\begin{prop}\label{prop:BoxTensorFunctorFullyFaithful}
The functor
$$
\boxtimes:
\Mod_G(X)\otimes\Mod_H(Y)
\hookrightarrow
\Mod_{G\times H}(X\times Y)
$$
% ---- 原文行 3260 ----
is fully faithful,
that is,
\begin{align*}
\boxtimes:
\Hom_{\Mod_G(X)}&
\big(
(\mathcal{F},\theta),
(\mathcal{F}',\theta')
\big)
\
\otimes
\
\Hom_{\Mod_H(Y)}
\big(
(\mathcal{G},\eta),
(\mathcal{G}',\eta')
\big)
\\ 
&\xrightarrow{\cong}
\Hom_{\Mod_{G\times H}(X\times Y)}
% ---- 原文行 3280 ----
\big(
(\mathcal{F},\theta)\boxtimes(\mathcal{G},\eta),
(\mathcal{F}',\theta')\boxtimes(\mathcal{G}',\eta')
\big)
\end{align*}
\end{prop}





